% Drawing 57/08-4 -- folder 57, batch 08, archive page 152 (PDF page 153).
% Transcription: une chaîne de composantes rationnelles marquée $g = 0$, avec « $(g = 0$ » ; un ovale ; un arc portant des points épaissis, accompagné de croix ; plusieurs étoiles de droites concourantes ; un pentagone de droites marqué de points ; deux lettres $k$ isolées
% Kind: page of sketches of stable curves. Method 2 (schematic redraw from the facsimile). Status: partial.
% Notes: Whole-page sketches, three groups. Top right: crossing strokes, and a chain of components - a long line crossed by three strokes ending in a blot - marked g = 0, with "(g = 0" further right. Lower left: a thick arc with two heavy end points crossed by a star of thin strokes; a second thick arc carrying three thickened points with small crosses beside it; the two letters k. Lower middle: an oval; a configuration of three lines marked 1; a pentagon of five lines carrying three points; two stars of concurrent lines and a small cross. The rays of the stars are approximate (number and direction read by eye).
% Coordinates: 1 unit = 1 mm = 1/100 of the crop width (bbox in page fractions: [0.13, 0.23, 0.92, 0.81]).
% Redrawn by Claude (claude-opus-5-5), issue #73, 2026-11-07; not checked by a human; not compiled (no TeX in the session).
% Rights: reproduces part of the work -- behind RENDER_TRANSCRIPTS, never published by itself.
\documentclass[tikz,border=3pt]{standalone}
\usepackage{amsmath,amssymb}
\usetikzlibrary{patterns,arrows.meta}
\tikzset{pencil/.style={black!45}, struck/.style={black!35}, mark/.style={circle,draw,inner sep=1.3pt}, pt/.style={circle,fill,inner sep=0.9pt}, opt/.style={circle,draw,fill=white,inner sep=1.1pt}}
\begin{document}
\begin{tikzpicture}[x=1mm,y=1mm,line cap=round,line join=round,font=\scriptsize]
\draw (65.5,-10.3) -- (75,-3.3);
\draw (73.9,-3.3) -- (77.6,-9.9);
\draw[very thick] (67,-5.1) -- (70.6,-14.7);
\draw (65.5,-18.4) -- (67.7,-12.1);
\draw (66.2,-12.8) -- (79.8,-16.5);
\node[pt] at (80.5,-16.2) {};
\draw (72.8,-12.1) -- (73.2,-19.5);
\draw (75,-12.1) -- (75,-19.1);
\draw (77.6,-13.6) -- (78,-20.6);
\node[inner sep=1pt] at (83.8,-16.9) {\tiny $g=0$};
\node[inner sep=1pt] at (95.6,-16.9) {\tiny $(g=0$};
\draw[very thick] plot[smooth, tension=0.55] coordinates {(9.3,-62.1) (11.1,-67.2) (14.2,-71.2) (18,-75.5)};
\node[pt] at (9.3,-62.1) {};
\node[pt] at (18,-75.5) {};
\draw[pencil] (7.1,-65.7) -- (13.8,-65.5);
\draw[pencil] (10.1,-67.6) -- (15.4,-63.5);
\draw[pencil] (8.1,-69.6) -- (17,-71.2);
\draw[pencil] (11.7,-72.4) -- (16.2,-68.6);
\draw[pencil] (12.6,-73.6) -- (14.6,-68.2);
\draw[very thick] plot[smooth, tension=0.55] coordinates {(5.1,-77.3) (5.1,-81.3) (6.7,-85) (9.1,-87.8) (12.8,-89.4)};
\node[pt] at (5.5,-81.9) {};
\node[pt] at (6.9,-85) {};
\node[pt] at (9.1,-87.8) {};
\node[inner sep=1pt] at (3.6,-82.8) {\tiny $\times$};
\node[inner sep=1pt] at (4.5,-85.8) {\tiny $\times$};
\node[inner sep=1pt] at (7.5,-89) {\tiny $\times$};
\node[inner sep=1pt] at (5.1,-95.1) {\small $k$};
\node[inner sep=1pt] at (14.2,-95.1) {\small $k$};
\draw (39.6,-61.4) ellipse[x radius=3, y radius=4.7];
\node[inner sep=1pt] at (23.3,-66.6) {\tiny $1$};
\draw (22.4,-74.4) -- (28.4,-65.7);
\draw (22.8,-68.8) -- (34.5,-70.5);
\draw (31,-66.6) -- (32.7,-75.6);
\draw (40.1,-79.5) -- (51.7,-70.9);
\draw (41.3,-78.2) -- (53,-74.4);
\draw (46.5,-74.8) -- (53.8,-77.4);
\draw (43.9,-79.5) -- (46.5,-74.8);
\draw (46.5,-74.8) -- (49.9,-79.1);
\node[pt] at (46.5,-74.8) {};
\draw (43.1,-70) -- (54.7,-69.2);
\draw (50.4,-65.3) -- (50.4,-73.1);
\draw (50.4,-69.6) -- (61.6,-64.5);
\draw (50.4,-69.6) -- (44.8,-79.9);
\draw (54.7,-73.1) -- (62.9,-70.9);
\draw (54.7,-73.1) -- (50.8,-77.4);
\draw (54.7,-73.1) -- (60.3,-78.7);
\draw (44.4,-85.1) -- (49.1,-81.2);
\draw (45.2,-82.1) -- (48.7,-85.5);
\draw (25.8,-83) -- (35.3,-84.3);
\draw (23.7,-85.5) -- (31.4,-89.4);
\draw (28.9,-81.7) -- (23.7,-89.8);
\draw (26.7,-90.7) -- (37.5,-86);
\draw (33.2,-83.8) -- (35.3,-90.3);
\node[pt] at (26.7,-84.7) {};
\node[pt] at (30.6,-83.4) {};
\node[pt] at (31.4,-88.1) {};
\end{tikzpicture}
\end{document}
